\section{Experiments}
\label{sec:setup}

\begin{table}[t]
\centering\scriptsize
\renewcommand{\arraystretch}{1.08}\setlength{\tabcolsep}{4pt}
\begin{tabular}{@{}p{12mm}p{17mm}p{18.5mm}p{19.5mm}@{}}
\toprule
& \textbf{W1} & \textbf{W2} & \textbf{W3} \\
\midrule
Tasks & 80 HotpotQA questions & 40 $\tau$-bench retail tasks & 80 FRAMES questions \\
Parts & one agent, live Wikipedia search & one agent, 16 shop tools, simulated customer &
planner, researchers, writer, reviewer, reviser \\
Variants & 18 (core set: 10) & 11 (core set: 8) & 13, in two groups \\
Clear pairs & 20 & 14 & 8 (planner), 5 (writer) \\
Ceiling & yes & no & yes \\
Full run, per task & \$0.003 & \$0.022 & \$0.006 \\
\bottomrule
\end{tabular}
\caption{The three workflows. We count the clear pairs on the variant set of the main results.}
\label{tab:workflows}
\end{table}


\paragraph{Workflows.}
Table~\ref{tab:workflows} shows the three workflows. They increase in size: one agent (W1), one
agent with a shop database and a simulated customer (W2), and a pipeline of five roles (W3). W1
answers multi-hop questions from HotpotQA \citep{yang2018hotpotqa} with live Wikipedia search. W2
uses the retail tasks of $\tau$-bench \citep{yao2024taubench}. W3 answers questions from FRAMES
\citep{krishna2024fact}. All parts of all workflows use one small model.

\paragraph{Variants.}
Each workflow has three types of variants. Edits are plausible improvements to a prompt, which a
separate LLM instance wrote. Mutants are variants that we broke on purpose, for example a limit of
two agent steps. Model swaps replace the main model with a stronger or a cheaper model. In W3, the
planner and the writer each have a group of variants.

\paragraph{Ground truth.}
Each variant ran once on each task under the full workflow. A second full run gave the noise
ceiling for W1 and for both groups of W3. For W2, the budget did not let us run a ceiling. Later
miniatures of W1 and W2 ran on a core set of variants. Each core set keeps all the mutants that the
full workflow catches.

\paragraph{Judge and scores.}
A blind LLM judge graded each run and gave one reason for each failure, from a fixed list for each
workflow. The judge did not know the variant or the recipe, or if a run came from a miniature. For
W2, the check of $\tau$-bench gives pass or fail. A simulated tool does not change the database.
Thus, we replay the calls of each run that change the database on a new copy of it, and we check
the result.

\paragraph{Cost.}
Appendix~\ref{app:record} gives the models, the number of runs and the cost. The full study cost
\$74.43 in API calls, under a hard limit that we checked before each call.
